\subsection{Why the degree tensor has the price of a shared label}\label{sec:degree}

Lemma~\ref{lem:vote} needs the price of \(W_m(C)\) for the resource it
supplies and the price of \(W_s(C)\) for the resource it returns. For
general \(s\), a direct degeneration extracts a smaller common label.
To obtain the exact price at powers of \(p\), we let the sites vary in
Euclidean space. The partial diagonals and the degree tensors have
the same price on a segment. A degeneration on intersecting sets of
sites will force their prices to agree on every convex hull.

For these comparisons use the joint Euclidean assignment of
\cite[Theorem~\ext{G}{thm:main} and Proposition~\ext{G}{prop:joint-convexity}]{oaG}.
Its unnormalized price for a nonzero tensor is jointly convex in
the placements of its factors, invariant under Euclidean isometries
and homogeneous under simultaneous dilation. At coincident sites the
factors form a single local space. Adding a scalar factor does not
change the price, and placing all factors together gives price zero.
On a fixed finite set \(S\) of positive width the unnormalized price
is \(D h\), where
\begin{equation}\label{eq:euclidean-price}
 D=\delta(S),\qquad
 w(P)=\frac{\delta(P)}{D},\qquad
 \delta(P)=\sqrt{2\pi}\,\mathbb E\max_{v\in P}\langle g,v\rangle ,
 \qquad g\sim N(0,I).
\end{equation}
The quantity \(\delta\) depends only on the convex hull, adds under
Minkowski sums and equals length on segments.

\begin{proof}[Proof of Lemma~\ref{lem:degree-width}]
Fix a power \(m>1\) of \(p\). For a nonempty finite set \(A\) in the
Euclidean flat, let \(f(A)\) be the unnormalized price of \(W_m(A)\)
divided by \(\log m\). The remainder construction of Section~\ref{sec:vote} gives a restriction
of \(G_A(m)\) at each \(z\ne0\) and gives \(W_m(A)\) at zero. Polynomial
specialization therefore gives \(0\le f(A)\le\delta(A)\). To prove equality we adapt the mean-width
argument for diagonal tensors in \cite[proof of Theorem~\ext{G}{thm:flatlaw}]{oaG}.

\emph{Intersecting sets of sites.}
For finite sets \(A,B\) with nonempty intersection \(I=A\cap B\), we
claim that
\begin{equation}\label{eq:degree-submod}
                 f(A\cup B)+f(I)\le f(A)+f(B).
\end{equation}
Represent \(W_m(A)\otimes W_m(B)\) by coefficient extraction at
\(t^{m-1}u^{m-1}\) in \(\mathbb F[t,u]/(t^m,u^m)\). The local
inputs are polynomials in \(t\) on \(A\setminus I\), in \(u\) on
\(B\setminus I\), and in both variables on \(I\). In the inputs from
\(B\setminus I\), substitute \(t+zu\) for \(u\).

At \(z\ne0\) this changes the form only by invertible local maps and a
nonzero scalar. Indeed, \(t\mapsto t,\ u\mapsto t+zu\) is a graded
automorphism of the quotient algebra because
\((t+zu)^m=t^m+z^m u^m\). Apply it to the full product after applying
its inverse to the inputs at \(I\). The automorphism scales the
one-dimensional top-degree component by \(z^{m-1}\).
At \(z=0\) every site contributes to the \(t\)-degree, and precisely
the sites in \(I\) contribute to the \(u\)-degree. The specialized
form is \(W_m(A\cup B)\otimes W_m(I)\). Taking prices proves
\eqref{eq:degree-submod}.

\emph{From segments to polytopes.}
Restricting the local degree at a new site to zero recovers the old
tensor, so \(f\) is monotone under inclusion. At a repeated site the
map \(v_e\mapsto\sum_{a+b=e} v_a\otimes v_b\) inserts an extra slot,
and selecting degree zero in that slot reverses the map. If a new
site lies in \(\operatorname{conv}(A)\), convexity bounds its price
by the price of adjoining it at one of the old sites. Monotonicity
gives the reverse bound. Hence \(f(A)\) depends only on
\(K=\operatorname{conv}(A)\).

The resulting function on polytopes satisfies
\(f(K+L)\le f(K)+f(L)\). Take vertex lists \((x_i)\) of \(K\) and
\((y_j)\) of \(L\), and give \(W_m\) one slot for each pair \((i,j)\).
For the placements \(x_i\) and \(y_j\) the values are \(f(K)\) and
\(f(L)\), since repeated sites have no effect. At the midpoint
placement the convex hull is \((K+L)/2\), so joint convexity and
homogeneity give the inequality.

For a vector \(v\) and an integer \(N\ge1\), form the chain of
translates \(A,A+v,\ldots,A+Nv\). When adjoining the adjacent pair
\((A+hv)\cup(A+(h+1)v)\) to the preceding union, the intersection
already contains \(A+hv\). By \eqref{eq:degree-submod} and translation
invariance, each increment in price is at most \(f(K+[0,v])-f(K)\).
The whole union contains \(x,x+Nv\) for any \(x\in A\). At two sites
\(W_m\) is a diagonal after reversal of one degree basis. Since its
unnormalized price divided by \(\log m\) is the distance, we obtain
\[
 N\|v\|\le f(K+[0,Nv])
              \le f(K)+N\bigl(f(K+[0,v])-f(K)\bigr).
\]
Dividing by \(N\) and taking \(N\to\infty\) gives
\(f(K+[0,v])\ge f(K)+\|v\|\). Subadditivity gives the other inequality.
Iterating from a point shows that \(f=\delta\) on every finite sum
of segments.

\emph{Averaging rotations.}
If two polytopes \(K,L\) have Hausdorff distance at most \(t\), each
lies in the sum of the other with \(tQ\), where \(Q\) is a fixed cube
containing the unit ball. Monotonicity and subadditivity give
\(\lvert f(K)-f(L)\rvert\le t\delta(Q)\). Thus \(f\) extends
continuously to compact convex sets. The Minkowski average of all
orthogonal images of \(K\) is a ball \(B\) centered at the origin,
since its support function is constant on the unit sphere.
Approximating the rotation average by finite averages gives
\(f(B)\le f(K)\) by subadditivity and invariance. The same averaging
applied to a segment expresses any origin-centered ball as a limit of
finite Minkowski sums of segments. Hence \(f(B)=\delta(B)\).
Invariance and Minkowski additivity of \(\delta\) give
\(\delta(B)=\delta(K)\), so
\[
                   \delta(K)=f(B)\le f(K)\le\delta(K).
\]
Normalizing as in \eqref{eq:euclidean-price} proves the power-of-\(p\)
formula. At \(m=1\) the tensor is the scalar unit.

\emph{Arbitrary remaining degree.}
Let \(|C|=a+1\ge2\). At \(a\) sites of \(W_s(C)\) retain degrees
\(x_1,\ldots,x_a\) in \(\{0,\ldots,\lfloor(s-1)/a\rfloor\}\).
At the final site retain degree \(s-1-a x_0\) with \(x_0\) in the
same interval. The remaining support satisfies \(\sum_{i=1}^a x_i=a x_0\).
Scale each term by \(z\) to the power
\[
      \sum_{i=1}^a x_i^2-a x_0^2=\sum_{i=1}^a(x_i-x_0)^2 .
\]
On the left the exponent is a sum of local integer weights, so the
scaling is local for \(z\ne0\). On the right it is nonnegative on
the support and vanishes precisely at equality. Specializing at zero
therefore gives a diagonal of size \(1+\lfloor(s-1)/a\rfloor
=\lceil s/a\rceil\), whose price is the claimed lower bound. For a
single site both the price and \(w(C)\) are zero.
\end{proof}
